\Needspace{8\baselineskip}\originalchapter{2003 作业五: 逆磁流与动量守恒}
题源: 14 Oct 2003, due 21 Oct.
\begin{exercise}\label{ex:03-5-1}
原题1. 无碰撞平衡, \ensuremath{\partial}y=\ensuremath{\partial}z=0、B=Bz均匀, n/T沿x变. 位于r粒子的y速度分布取中心r0的Maxwell: $f=n(r_0)[m/(2\pi T(r_0))]^{1/2}\exp[-mv_y^2/(2T(r_0))]$, 且 $\mathbf r_0-\mathbf r=-\mathbf B\times m\mathbf v_\perp/(qB^2)$. 一阶回旋半径展开, 求粒子局域平均vy并核对流体逆磁速度.
\end{exercise}
\begin{solution}
$r_{0x}=x+mv_y/(qB)$, 故 $f=f_M(x,v_y)+[mv_y/(qB)]\partial_xf_M$. 零阶vy积分零, 一阶 \begin{equation}n\langle v_y\rangle=\frac m{qB}\partial_x\int v_y^2 f_M\dd v_y=\frac1{qB}\partial_x(nT).\end{equation} 因而 $\mathbf v_*=(\mathbf B\times\nabla p)/(qnB^2)=p'\mathbf e_y/(qnB)$. 它是位置分布偏差产生的局域粒子流, 不意味着中心集体穿越压强面; 与上一套中心电流/磁化区别一致.
\end{solution}
\begin{exercise}\label{ex:03-5-2}
原题2. 证明单流体加电磁场动量可写 $\partial_t\mathbf G=-\nabla\cdot\boldsymbol\Pi$, $\mathbf G=\rho\mathbf V+\epsilon_0\mathbf E\times\mathbf B$, $\boldsymbol\Pi=\rho\mathbf V\mathbf V+pI-[\mathbf B\mathbf B/\mu_0+\epsilon_0\mathbf E\mathbf E-(B^2/(2\mu_0)+\epsilon_0E^2/2)I]$.
\end{exercise}
\begin{solution}
连续性将物质动量写 $\partial_t(\rho V_i)+\partial_j(\rho V_iV_j+p\delta_{ij})=\rho_qE_i+(\mathbf J\times\mathbf B)_i$. Maxwell含位移电流给Lorentz力 $\rho_q\mathbf E+\mathbf J\times\mathbf B=\nabla\cdot\mathbf T_M-\partial_t(\epsilon_0\mathbf E\times\mathbf B)$. 移项即所给\ensuremath{\Pi}, 负Maxwell应力是\textbf{动量通量}而非受力应力. 若使用严格低速准中性MHD并删位移电流, E能和场动量为更高阶小量, 上式相应按同阶近似理解; 精确包含E项时必须使用完整Maxwell, 不混用两种阶数.
\end{solution}
\begin{exercise}\label{ex:03-5-3}
原题3. 氢tokamak Te=Ti=1keV,n=$5\times10^{19}$m$^{-3}$, 密度梯度$\dd n/\dd r=n/L$, $L=.2$ m.
 B\ensuremath{\phi}=2T,B\ensuremath{\theta}=.2T,J=$10^6$A/m\ensuremath{{}^{2}}, 每量dg/dt=g/τ、τ=1s. (a)J纯环向、V=0, 一位有效数字估全Ohm各项; (b)J方向改为满足\ensuremath{\nabla}p=J\ensuremath{\times}B, 重算改变的项.
\end{exercise}
\begin{solution}
用式\eqref{eq:ohm-general}, 取ln\ensuremath{\Lambda}\ensuremath{\simeq}17. (a)\ensuremath{\eta}J约 $3\times10^{-2}$V/m; Hall $JB_\theta/(en)=2.5\times10^4$V/m; 电子压力 $T_e/(eL)=5\times10^3$V/m; 电子惯性 $m_eJ/(ne^2\tau)=7\times10^{-7}$V/m; V\ensuremath{\times}B=0. E不是独立已知值, 由这些向量项求出, 径向项需按方向相减而非只相加模. 因未给曲率/径向流, 对流惯性不能从τ凭空指定; 时间惯性为上列估计.

(b)总 $p=p_e+p_i=2nT$, 故 Hall $|\nabla p|/(en)=2T/(eL)=10^4$V/m, 压力仍 $5\times10^3$V/m, 径向E\ensuremath{\simeq}$5\times10^3$V/m. 电流垂直B分量约 $|\nabla p|/B=4\times10^4$A/m\ensuremath{{}^{2}}, 其余主要平行, 总电阻与时间惯性量级不变. 纯环向假设(a)不满足给定压强平衡, 不应把两情形混称同一平衡.
\end{solution}
